\section{When a batch goes}
\label{sec:ch04-when-a-batch-goes}
\index{DataLoader!batching internals|(}

A DataLoader collects keys and then, at some point, stops collecting and
fetches. That word, \enquote{then}, is the whole of the mechanism. Stop too
early and you have several small batches; stop too late and every request pays
for the wait. The documentation describes the trigger like this:

\begin{quote}
\enquote{once no more resolver work is immediately ready, each pending batch is
dispatched \dots\ The dispatch trigger is \enquote{no more ready work}, not a
fixed schedule and not one dispatch per level of the
query.}~\autocite{chillicream2026dataloaderdocs}
\end{quote}

The second sentence is right and the first is not the mechanism. Nothing in the
code that dispatches batches observes resolver readiness. It observes the batch.
This chapter's rule, set in chapter~\ref{ch:the-life-of-a-request}, is that where
the source and a page disagree the source wins, so here is the source.

\subsection*{The half in GreenDonut: collecting}

\index{GreenDonut!DataLoaderBase}%
\code{DataLoaderBase<TKey,TValue>} keeps one open batch. \code{LoadAsync} takes
a lock, asks the promise cache for the key, and on a miss adds it to that batch:

\begin{minted}{csharp}
var current = _currentBatch;

if (current is not null)
{
    // if the batch has space for more keys we just keep adding to it.
    if (current.Size < _maxBatchSize || _maxBatchSize == 0)
    {
        return current.GetOrCreatePromise<TValue?>(key, allowCachePropagation);
    }
    ...
}
\end{minted}

\code{MaxBatchSize} defaults to 1024. Mosaic's largest batch is twenty-five keys,
so the branch below that comment has never run here, and I have not measured what
happens when it does.

\index{GreenDonut!Batch}%
What the batch records is the interesting part, because it is what the other half
reads. Every \code{Batch} carries a creation timestamp, a modification timestamp,
a status of either enqueued or touched, and this:

\begin{minted}{csharp}
public override bool Touch()
{
    var previous = Interlocked.Exchange(ref _status, Touched);
    return previous == Touched;
}
\end{minted}

Adding a key sets the status back to enqueued and restamps the modification
time. So \code{Touch} returning true means precisely one thing: nothing has put a
key in this batch since the last time somebody looked.

\subsection*{The half in HotChocolate: deciding}

\index{HotChocolate!BatchDispatcher}%
\code{BatchDispatcher} is registered scoped, so there is one per request. Its
\code{Schedule} method does no work at all; it adds the batch to a set,
increments a counter, and returns. The work is in a coordinator loop that runs
for the life of the request, and its decision is two lines:

\begin{minted}{csharp}
var shouldDispatch =
    batch.Touch()
    && TicksToUs(now - batch.ModifiedTimestamp) >= batchSettleTimeUs;

if (!shouldDispatch && maxBatchAgeUs != 0)
{
    shouldDispatch = TicksToUs(now - batch.CreatedTimestamp) > maxBatchAgeUs;
}
\end{minted}

\index{HotChocolate!BatchDispatcherOptions}%
Two conditions, not one. A batch goes when it has survived two evaluation rounds
without anybody touching it \emph{and} has been quiet for at least the settle
time, or when it is simply too old. The defaults are 250 microseconds for the
settle time and 50 milliseconds for the ceiling, and the ceiling can be turned
off by setting it to zero.

Between rounds the coordinator yields, and the comment on that yield is the
clearest statement of the design in the file:

\begin{minted}[fontsize=\footnotesize]{csharp}
// Yield between evaluation rounds so executing resolvers get time
// to add more data requirements to the open batches. Under load the
// yield queues behind real work, which widens the aggregation window.
\end{minted}

That is theirs, not mine. It also names a property I have not measured and will
not claim: under load the window widens by itself, because the coordinator's
yield goes to the back of a busier queue. It is a nice piece of design and this
chapter's evidence is one client on one machine.

\index{HotChocolate!batch priority}%
When several batches are open they go into a priority queue keyed on the
modification timestamp, so the batch nobody has touched for longest is evaluated
first and has the best chance of being dispatched. A single open batch skips the
queue entirely through a separate method. That ordering is why a query with
several independent DataLoaders does not starve one of them behind another.

\subsection*{The scheduler you get when nobody is coordinating}

\index{GreenDonut!AutoBatchScheduler}%
There is a second implementation, and knowing it exists explains a class of
confusing benchmark. \code{AutoBatchScheduler} is what a DataLoader gets outside
a HotChocolate request, and its \code{Schedule} is a single expression:

\begin{minted}{csharp}
public void Schedule(Batch batch)
    => BeginDispatch(batch);

private static void BeginDispatch(Batch batch)
    => Task.Run(async () =>
    {
        try
        {
            await batch.DispatchAsync().ConfigureAwait(false);
        }
        catch
        {
            // we do ignore any potential exceptions here
        }
    });
\end{minted}

Its own summary says it \enquote{prioritizes low latency over batching
efficiency by executing each batch as soon as it is scheduled}. No settle time,
no coordination. Inside a request the scoped dispatcher is substituted in its
place, which happens through a small service provider wrapper that intercepts
requests for \code{IBatchScheduler}. If you ever test a DataLoader outside the
executor and find it batching far less than you expected, this is why, and the
answer is not that your DataLoader is broken.

\subsection*{The number is three, nearly always}

\index{DataLoader!split batches}%
Two conditions and a timer mean the answer is a distribution rather than a
constant, and it is worth knowing the shape of it rather than assuming.

I sent the catalog-with-reviews query four hundred times against a warm process
and counted what the timeline reported. Three hundred and ninety-eight requests
reported three statements. Two reported four. In both of those, the duplicated
statement was the one against \code{customers}, issued twice with different key
sets.

\index{DataLoader!fan-out}%
Which is exactly the mechanism above, seen from outside. The customers batch is
fed by a hundred and twenty resolvers, more than any other batch in the request,
and they do not all arrive in an instant. When the stragglers are late by more
than a quarter of a millisecond, the coordinator has seen the batch untouched
twice, the settle time has elapsed, and it dispatches with what it has. The
remaining keys open a new batch. Both statements are correct, the response is
identical, and it took two runs in four hundred to see it happen.
Figure~\ref{fig:ch04-batch-window} puts the two cases side by side.

\begin{figure}[htbp]
  \centering
  % Why a batch usually goes once and occasionally twice. Referenced from
% chapter 04. Bare tikzpicture: the figure environment, caption and label live
% at the call site. Uses only arrows.meta and positioning.
\begin{tikzpicture}[
    font=\small,
    axis/.style={-{Stealth[length=2mm]}, thick},
    key/.style={draw, fill=black!25, circle, inner sep=0pt, minimum size=1.6mm},
    late/.style={draw, fill=white, circle, inner sep=0pt, minimum size=1.6mm},
    win/.style={draw, dashed, rounded corners=2pt, minimum height=6mm},
    go/.style={-{Stealth[length=2mm]}, thick},
    lbl/.style={font=\scriptsize, align=left},
    note/.style={font=\scriptsize\itshape, align=left},
    head/.style={font=\scriptsize\bfseries}
  ]

  % ---- upper: everything inside the window -----------------------------
  \node[head, anchor=west] at (-0.4,1.55) {one batch};

  \draw[axis] (-0.4,0.8) -- (9.4,0.8);
  \foreach \x in {0.2,0.5,0.8,1.1,1.4,1.7,2.0,2.3}{%
    \node[key] at (\x,0.8) {};
  }
  \node[win, minimum width=17mm] at (3.2,0.8) {};
  \node[lbl, anchor=south] at (3.2,1.05) {quiet 250\,$\mu$s};
  \draw[go] (4.1,0.8) -- (5.3,0.8);
  \node[lbl, anchor=west] at (5.4,0.8) {dispatch: 8 keys, 1 statement};

  \node[note, anchor=west] at (-0.4,0.25)
    {every key arrives, then the coordinator sees the batch untouched twice};

  % ---- lower: a straggler misses it ------------------------------------
  \node[head, anchor=west] at (-0.4,-0.85) {two batches};

  \draw[axis] (-0.4,-1.6) -- (9.4,-1.6);
  \foreach \x in {0.2,0.5,0.8,1.1,1.4,1.7}{%
    \node[key] at (\x,-1.6) {};
  }
  \node[win, minimum width=17mm] at (2.6,-1.6) {};
  \node[lbl, anchor=south] at (2.6,-1.35) {quiet 250\,$\mu$s};
  \draw[go] (3.5,-1.6) -- (4.3,-1.6);
  \node[lbl, anchor=west] at (4.4,-1.6) {dispatch: 6 keys};

  \node[late] at (5.0,-2.3) {};
  \node[late] at (5.3,-2.3) {};
  \draw[axis, dashed] (4.7,-2.3) -- (9.4,-2.3);
  \node[win, minimum width=17mm] at (6.4,-2.3) {};
  \draw[go] (7.3,-2.3) -- (8.0,-2.3);
  \node[lbl, anchor=west] at (8.1,-2.3) {2 keys};

  \node[note, anchor=west] at (-0.4,-3.1)
    {two stragglers arrive after the batch has gone, and open a new one};

\end{tikzpicture}

  \caption{Why the answer is usually three and occasionally four. The
  coordinator dispatches a batch it has seen untouched across two rounds and
  quiet for the settle time. On the upper timeline every key arrives inside the
  window. On the lower one a straggler misses it, so the batch goes with what it
  has and the rest form a second one. Nothing is wrong in either case.}
  \label{fig:ch04-batch-window}
\end{figure}

This had a consequence for the companion repository's verification script, and
the way it was resolved is the part I would defend. The script asserts the
numbers the book prints, so it wanted to assert three. Asserting three against a
single request would fail roughly one run in two hundred, and a gate that fails
at random gets ignored within a week. Accepting \enquote{three or four} from one
request would never fail, and would also stop distinguishing the ordinary case
from the exception, which is most of what the assertion is for.

So it sends the query five times and asserts two things: that at least one run
reported exactly three, and that no run reported more than four. A service
without DataLoaders reports a hundred and forty-six five times over and fails
both. That is the general shape of the problem, and it is not specific to this
book: a number produced by a race needs an assertion that describes the race,
and the way to write one is to measure the distribution first.

Three statements is where the N+1 story ends. The next section is about what is
in them.
\index{DataLoader!batching internals|)}
